\documentclass[10pt,a4paper]{amsart}
\usepackage[margin=19mm]{geometry}
\usepackage{amsmath,amssymb,mathtools}
\usepackage[scr=rsfs,cal=euler]{mathalfa}
\usepackage{xfrac}
\usepackage[unicode,hidelinks,bookmarks=false]{hyperref}
\newcommand{\ie}{i.e.\ }
\newcommand{\cf}{cf.\ }
\newcommand{\resp}{respectively\ }
\newcommand{\Subsection}{\subsection}
\providecommand{\nobreakdash}{\nobreak\hbox{-}}
\setlength{\emergencystretch}{2em}
\multlinegap0pt
\usepackage{amssymb}
\usepackage{enumerate}
\usepackage{tikz}
\usepackage{float}
\newtheorem{lemma}{Lemma}
\title{Optimal bounds for embedded eigenvalues of one-dimensional discrete Schr\"odinger operators with decaying potentials}
\author{Jifeng Chu, Wencai Liu, Kang Lyu}
\date{Source: arXiv:2609.03224v1 (CC BY 4.0)}
\begin{document}
\maketitle

\noindent\emph{Source and licence.} Optimal bounds for embedded eigenvalues of one-dimensional discrete Schr\"odinger operators with decaying potentials. Version arXiv:2609.03224v1 (CC BY 4.0), distributed by arXiv under the Creative Commons Attribution 4.0 International licence, http://creativecommons.org/licenses/by/4.0/, as recorded in the arXiv licence metadata for this exact version and shown on the abstract page https://arxiv.org/abs/2609.03224v1. Full licence notice: \texttt{licenses/arxiv-2609.03224-CC-BY-4.0.txt}.\par\smallskip

\noindent\textit{Editorial scope.} Counted unit: the article's lemma lssccs (source0.tex lines 831--848). The article's complete original proof of it is delivered with it (source0.tex lines 849--895, one \texttt{proof} environment of 47 lines). No mathematics is written, repaired or re-derived: the statement and the proof are the article's own lines, in the article's own notation, and no retained line is substituted or re-flowed.  Retained uncounted as the source-local context the proof needs: lines 765--767; lines 773--775; lines 808--810.  Nothing was omitted from any retained span, so the package carries no omission record.  No retained line contains a citation, so the wrapper carries no bibliography and no bibliography entry.  No retained line contains a figure environment, an image-inclusion command or drawing code, so no image is delivered and the package declares no asset dependency.  Editorial formatting only, in addition: an AMS \texttt{amsart} preamble that restates the article's own declarations and macros used by the retained lines, namely the environments \texttt{lemma} on separate counters, together with the standard packages those lines need.  The retained environments print in this document as Lemma 1 in order of appearance, whereas the article prints the same items with the numbering of its own full document.  The complete native source of the article as distributed in the arXiv e-print for this exact version is bundled unchanged as \texttt{sources/209279/source0.tex}.
\begin{align}
	T(\varphi)=\sum_{j=0}^{q-1}\sin^2\pi\left(\varphi+\frac{j}{q}\right){\rm{sgn}}\left(\sin\left(2\pi\varphi+\frac{2\pi}{q}j\right)\right).\label{definitionofT_qtheta}
\end{align}

\begin{align}
		T(\varphi)=T\left(\varphi+\frac{1}{q}\right)\label{perioTvarphi}.
\end{align}

		\begin{align}
			\sum_{i=0}^n\sin (\varphi+ix)=\frac{\cos\left(\varphi-\frac{x}{2}\right)-\cos\left(\varphi+\left(n+\frac{1}{2}\right)x\right)}{2\sin\frac{x}{2}}.\label{sintheta+ix}
		\end{align}

	\begin{lemma}\label{lssccs}Assume that $q\geq3$ is odd. Then for any $\varphi\in \mathbb{R}$, one has
		\begin{align}\label{tqduichen}
			T(\varphi)=-T\left(\frac{1}{q}-\varphi\right).
		\end{align}
		Specifically,	we have
		\begin{align}
			T(\varphi)=\begin{cases}
				\label{ssccs}\frac{1}{2}+\frac{\sin2\pi\varphi-\sin\left(\frac{\pi}{q}-2\pi\varphi\right)}{2\sin\frac{\pi}{q}}, & \varphi\in[0,\frac{1}{2q}),\\
				0, & \varphi=\frac{1}{2q},\\
				-\frac{1}{2}-\frac{\sin\left(\frac{2\pi}{q}-2\pi\varphi\right)+\sin\left(\frac{\pi}{q}-2\pi\varphi\right)}{2\sin\frac{\pi}{q}}, & \varphi\in  (\frac{1}{2q}, \frac{1}{q}],
			\end{cases}
		\end{align}
		and
		\begin{align}
			\label{mp1}&\lim_{\varphi\to\frac{1}{2q}\pm0}T(\varphi)=\mp1,\\
			\label{tqpdy0}T^{\prime}(\varphi)>0&, \mathrm{\ for\ any\ }\varphi\in\left(0,\frac{1}{2q}\right)\cup\left(\frac{1}{2q},\frac{1}{q}\right).
		\end{align}
	\end{lemma}

	\begin{proof}
		Once we have \eqref{ssccs}, direct computations will show \eqref{mp1}, \eqref{tqpdy0}, and \eqref{tqduichen} for any $\varphi\in\left[0,\frac{1}{q}\right]$. Then by \eqref{perioTvarphi} one can obtain \eqref{tqduichen} for any $\varphi\in\mathbb{R}$.
		We only prove \eqref{ssccs}.
		
		Applying \eqref{sintheta+ix}, one can obtain that for any $\varphi\in\mathbb{R}$,
		\begin{align}
			\sum_{i=0}^n\cos(\varphi+ix)=\frac{\sin\left(\varphi+\left(n+\frac{1}{2}\right)x\right)-\sin\left(\varphi-\frac{x}{2}\right)}{2\sin\frac{x}{2}}.\label{sumcostheta+ix}
		\end{align}
		Hence, one has
		\begin{align}
			\sum_{j=0}^{q-1}\cos\left(2\pi\varphi+\frac{2\pi}{q}j\right)=0.\label{sumj=0q-1cos=0}
		\end{align}
	By \eqref{definitionofT_qtheta}, for any $\varphi\in\Big[0,\frac{1}{2q}\Big)$, 
		\begin{align}
			T(\varphi)
			=&\sum_{j=0}^{\frac{q-1}{2}}\sin^2\pi\left(\frac{j}{q}+\varphi\right )-\sum_{j=\frac{q+1}{2}}^{q-1}\sin^2\pi\left(\frac{j}{q}+\varphi\right)\nonumber\\
			=&\sum_{j=0}^{\frac{q-1}{2}}\frac{1-\cos\left(2\pi\varphi+\frac{2\pi}{q}j\right)}{2}-\sum_{j=\frac{q+1}{2}}^{{q-1}}\frac{1-\cos\left(2\pi\varphi+\frac{2\pi}{q}j\right)}{2}.\nonumber
		\end{align}
Using \eqref{sumcostheta+ix} and \eqref{sumj=0q-1cos=0}, one has
\begin{align}
	T(\varphi)=\frac{1}{2}-\sum_{j=0}^{\frac{q-1}{2}}\cos\left(2\pi\varphi+\frac{2\pi}{q}j\right)=\frac{1}{2}+\frac{\sin2\pi\varphi-\sin\left(\frac{\pi}{q}-2\pi\varphi\right)}{2\sin\frac{\pi}{q}}.\nonumber
\end{align}
Thus, we obtain \eqref{ssccs} for  $\varphi\in [0,\frac{1}{2q})$.

By \eqref{definitionofT_qtheta},  for $\varphi=\frac{1}{2q}$,
\begin{align}
	T\left(\frac{1}{2q}\right)
	=&\sum_{j=0}^{\frac{q-3}{2}}\frac{1-\cos\left(\frac{\pi}{q}+\frac{2\pi}{q}j\right)}{2}-\sum_{j=\frac{q+1}{2}}^{{q-1}}\frac{1-\cos\left(\frac{\pi}{q}+\frac{2\pi}{q}j\right)}{2}\nonumber\\
	=&-\frac{1}{2}\left(\sum_{j=0}^{\frac{q-3}{2}}\cos\left(\frac{\pi}{q}+\frac{2\pi}{q}j\right)-\sum_{j=\frac{q+1}{2}}^{{q-1}}\cos\left(\frac{\pi}{q}+\frac{2\pi}{q}j\right)\right).\nonumber
\end{align}
Since for any $j\in\left\{0,1,\cdots,\frac{q-3}{2}\right\}$ and any $l\in\left\{\frac{q+1}{2},\cdots,q-1\right\}$ with $j+l=q-1$, 
\begin{align}
\cos\left(\frac{\pi}{q}+\frac{2\pi}{q}j\right)=\cos\left(\frac{\pi}{q}+\frac{2\pi}{q}l\right),\nonumber
\end{align}
we have $T\left(\frac{1}{2q}\right)=0$.

		 By \eqref{definitionofT_qtheta}, for any $\varphi\in \Big(\frac{1}{2q},\frac{1}{q}\Big]$,
		\begin{align}
			T(\varphi)=&\sum_{j=0}^{\frac{q-3}{2}}\sin^2\pi\left(\frac{j}{q}+\varphi\right )-\sum_{j=\frac{q-1}{2}}^{q-1}\sin^2\pi\left(\frac{j}{q}+\varphi\right)\nonumber\\
			=&\sum_{j=0}^{\frac{q-3}{2}}\frac{1-\cos\left(2\pi\varphi+\frac{2\pi}{q}j\right)}{2}-\sum_{j=\frac{q-1}{2}}^{{q-1}}\frac{1-\cos\left(2\pi\varphi+\frac{2\pi}{q}j\right)}{2}\nonumber.
		\end{align}
		Hence, applying \eqref{sumcostheta+ix} and \eqref{sumj=0q-1cos=0}, one can obtain 
		\begin{align}
			T(\varphi)=-\frac{1}{2}-\sum_{j=0}^{\frac{q-3}{2}}\cos\left(2\pi\varphi+\frac{2\pi}{q}j\right)=-\frac{1}{2}-\frac{\sin\left(\frac{2\pi}{q}-2\pi\varphi\right)+\sin\left(\frac{\pi}{q}-2\pi\varphi\right)}{2\sin\frac{\pi}{q}}.\nonumber
		\end{align}
		Thus, we obtain \eqref{ssccs} for $\varphi\in\Big(\frac{1}{2q},\frac{1}{q}\Big].$
	\end{proof}

\end{document}
